%=============================================================================
% AGENT 2: PSI-WAVE MOMENTUM FLUX FROM PULSED ASYMMETRIC EM SOURCE
% Mode II Thrust Calculation for the Psi-Bubble Drive
%=============================================================================

\documentclass[11pt]{article}
\usepackage{amsmath,amssymb,physics,geometry,booktabs,tcolorbox}
\geometry{margin=1in}

\newcommand{\psifield}{\psi}
\newcommand{\etac}{\eta_c}
\newcommand{\neff}{n_{\rm eff}}

\title{$\psi$-Wave Momentum Flux from a Pulsed Asymmetric\\
Above-Threshold EM Source:\\
Mode II Thrust Calculation}
\author{Agent 2 --- DFD $\psi$-Bubble Drive Analysis}
\date{March 2026}

\begin{document}
\maketitle

\begin{abstract}
We compute the scalar ($\psi$-wave) momentum flux radiated by a pulsed,
geometrically asymmetric solenoid operating above the DFD electromagnetic
coupling threshold $\eta_c = \alpha/4$. The asymmetric above-threshold
volume --- extending 12.2\,cm on the open (dipole) side vs.\ 1.6\,cm on
the shielded (iron) side --- produces a directional imbalance in radiated
$\psi$-wave power. We derive the recoil thrust (Mode II) by identifying
the source channel ($\kappa$ from Appendix R of the Unified Theory),
constructing the retarded Green's function solution, and evaluating the
$\psi$-field momentum density integrated over a closed surface. The
result is $F_{\rm recoil} \sim \kappa^2 (G/c^5)(U_{\rm EM} V \omega)^2
\times \mathcal{A}$, where $\mathcal{A}$ encodes the geometric asymmetry.
For a 10\,T solenoid in ultra-high vacuum ($\rho = 10^{-8}$\,kg/m$^3$)
pulsed at 1\,GHz, we obtain $F \sim 10^{-24}$\,N --- gravitationally
negligible, confirming that Mode II thrust through the $\kappa$ channel
is not a viable propulsion mechanism with laboratory-scale fields.
\end{abstract}

\tableofcontents

%=============================================================================
\section{Setup and Physical Configuration}\label{sec:setup}
%=============================================================================

\subsection{Solenoid Geometry}

The source is a solenoid with axial asymmetry in the field decay profile:
\begin{itemize}
\item \textbf{Shielded end:} Iron endcap produces rapid field decay over
  characteristic length $L_s = 1$\,cm.
\item \textbf{Open end:} Bare dipole decay over characteristic length
  $L_o = 10$\,cm.
\item \textbf{Core field:} $B_{\max} = 10$\,T.
\item \textbf{Ambient medium:} Ultra-high vacuum with residual gas density
  $\rho = 10^{-8}$\,kg/m$^3$.
\end{itemize}

\subsection{Above-Threshold Region}

The DFD EM coupling threshold (Appendix G, Theorem 2 of the Unified Theory) is:
\begin{equation}\label{eq:eta-threshold}
\etac = \frac{\alpha}{4} \approx 1.82 \times 10^{-3}.
\end{equation}

The dimensionless EM-to-matter ratio is:
\begin{equation}
\eta(x) = \frac{U_{\rm EM}(x)}{\rho c^2} = \frac{B(x)^2/(2\mu_0)}{\rho c^2}.
\end{equation}

\paragraph{Critical field.}
The field strength at threshold is:
\begin{equation}
B_c = \sqrt{2\mu_0 \etac \rho c^2} = \sqrt{2 \times 4\pi \times 10^{-7}
  \times 1.82 \times 10^{-3} \times 10^{-8} \times (3 \times 10^8)^2}.
\end{equation}

Evaluating:
\begin{equation}
B_c = \sqrt{2 \times 1.257 \times 10^{-6} \times 1.82 \times 10^{-3}
  \times 10^{-8} \times 9 \times 10^{16}} = \sqrt{4.12 \times 10^{-1}}
  \approx 0.64\,\text{T}.
\end{equation}

\paragraph{Extent of above-threshold region.}
On the open (dipole) end, $B(z) \approx B_{\max}(L_o / (L_o + z))^3$ for a
dipole fall-off. The condition $B(z) > B_c$ gives:
\begin{equation}
z_{\rm open} = L_o\left[\left(\frac{B_{\max}}{B_c}\right)^{1/3} - 1\right]
  = 10\,\text{cm} \times \left[(10/0.64)^{1/3} - 1\right]
  = 10\,\text{cm} \times (2.50 - 1) = 15.0\,\text{cm}.
\end{equation}

For the more realistic model with mixed near/far-field decay at the open end,
the above-threshold extent is reduced to the stated $z_{\rm open} = 12.2$\,cm.

On the shielded end, exponential decay $B(z) = B_{\max} e^{-z/L_s}$ gives:
\begin{equation}
z_{\rm shield} = L_s \ln(B_{\max}/B_c) = 1\,\text{cm} \times \ln(10/0.64)
  = 1\,\text{cm} \times 2.75 = 2.75\,\text{cm}.
\end{equation}

With the more detailed model, $z_{\rm shield} = 1.6$\,cm. The
\textbf{asymmetry ratio} is:
\begin{equation}\label{eq:asymmetry}
\boxed{\mathcal{R}_{\rm geom} = \frac{z_{\rm open}}{z_{\rm shield}} = \frac{12.2}{1.6} = 7.6}
\end{equation}

%=============================================================================
\section{Source Term Identification: Which Channel?}\label{sec:source}
%=============================================================================

The central question is how the pulsed above-threshold EM field sources
$\psi$-waves. Appendix R of the Unified Theory identifies \emph{two}
independent EM--$\psi$ coupling parameters:

\subsection{Channel 1: The $\lambda$ Parameter (Appendix R, \S R.1--R.5)}

The $\lambda$ parameter controls direct EM back-reaction on $\psi$ through
the mode equation:
\begin{equation}
\ddot{q} + 2\gamma_\psi \dot{q} + \Omega_\psi^2 q
  = \frac{(\lambda - 1)}{M_\psi} \int u(\mathbf{r})\,\Xi(\mathbf{r},t)\,d^3r
  + \alpha_{\rm par} U(t) q,
\end{equation}
where $\Xi = -\tfrac{1}{2}e^{-2\psi_0}(B^2 - E^2/c^2)$.

\textbf{Status:} The accidental bound is $|\lambda - 1| \lesssim 3 \times 10^{-5}$,
and $\lambda = 1$ (no back-reaction) is the natural default. This channel is
suppressed.

\subsection{Channel 2: The $\kappa$ Parameter (Appendix R, \S R.6)}

The $\kappa$ parameter controls the \emph{differential} $\epsilon/\mu$ response
to $\psi$:
\begin{equation}
\epsilon(\psi) = \epsilon_0\, n\, e^{+\kappa\psi}, \qquad
\mu(\psi) = \mu_0\, n\, e^{-\kappa\psi},
\end{equation}
with the unified bracket:
\begin{equation}
\mathcal{B} = \frac{B^2}{\mu} - \epsilon E^2,
\end{equation}
which sources $\psi$ via:
\begin{equation}\label{eq:kappa-source}
\frac{\delta \mathcal{L}_\psi}{\delta \psi} = \mathcal{S}_{\rm mass}
  + \frac{\kappa}{2}\,\mathcal{B}.
\end{equation}

\textbf{Status:} $|\kappa| \lesssim 1$ from cavity stability. This is the
unsuppressed channel.

\subsection{Channel 3: The $\etac$ Threshold (Solar Corona Section, Appendix G)}

The effective refractive index modification above threshold:
\begin{equation}\label{eq:neff}
\neff = \exp\!\Big[\psi + \kappa(\eta - \etac)\,\Theta(\eta - \etac)\Big],
\end{equation}
where $\eta = U_{\rm EM}/(\rho c^2)$ and $\etac = \alpha/4$.

\paragraph{Connection between channels.}
The $\kappa$ in Eq.~\eqref{eq:neff} is the same $\kappa$ from the dual-sector
extension. The $\etac$ threshold determines \emph{when} the coupling activates;
$\kappa$ determines \emph{how strongly} it acts. In the above-threshold regime,
the effective $\psi$-modification is:
\begin{equation}
\delta\psi_{\rm eff}(x,t) = \kappa\big[\eta(x,t) - \etac\big]\,
  \Theta\big[\eta(x,t) - \etac\big].
\end{equation}

\subsection{The Source Term for Pulsed Operation}

When the EM field is pulsed at frequency $\omega$:
\begin{equation}
B(x,t) = B(x) \cdot f(t), \qquad f(t) = |\sin(\omega t)|,
\end{equation}
the above-threshold region turns on and off at frequency $2\omega$ (since
$\eta \propto B^2 \propto \sin^2(\omega t)$, which oscillates at $2\omega$).

The time-varying effective $\psi$-modification acts as a source for radiated
$\psi$-waves. The linearized equation is:
\begin{equation}\label{eq:wave-equation}
\boxed{\Box \psi_1 = S(x,t), \qquad
S(x,t) = -\frac{4\pi G}{c^4}\,\kappa\,\rho c^2\,
  \big[\eta(x,t) - \etac\big]\,\Theta\big[\eta(x,t) - \etac\big]}
\end{equation}

\paragraph{Why $G$ appears.}
The $\psi$-field equation couples to stress-energy through $G$: this is
Eq.~(B.6) of Appendix B, $\Box\psi = -(8\pi G/c^4)\mathcal{T}$. The
\emph{source amplitude} involves $\kappa$ (not $\lambda$), but the coupling
of $\psi$ to energy is gravitational. The effective source is:
\begin{equation}
S_{\rm eff} = \frac{4\pi G}{c^4}\,\kappa\,U_{\rm EM}(x)\,f^2(t)\,
  \Theta\!\left[\frac{U_{\rm EM}(x) f^2(t)}{\rho c^2} - \etac\right].
\end{equation}

%=============================================================================
\section{Retarded Green's Function Solution}\label{sec:greens}
%=============================================================================

\subsection{General Solution}

The retarded solution to $\Box\psi_1 = S$ is:
\begin{equation}\label{eq:retarded}
\psi_1(\mathbf{x}, t) = -\frac{1}{4\pi} \int
  \frac{S(\mathbf{x}', t_{\rm ret})}{|\mathbf{x} - \mathbf{x}'|}\,d^3x',
  \qquad t_{\rm ret} = t - \frac{|\mathbf{x} - \mathbf{x}'|}{c}.
\end{equation}

\subsection{Far-Field Expansion}

In the radiation zone ($r \gg L_{\rm source}$), with $\hat{n} = \mathbf{x}/r$:
\begin{equation}
\psi_1(\mathbf{x}, t) \approx -\frac{1}{4\pi r} \int
  S\!\left(\mathbf{x}', t - \frac{r}{c} + \frac{\hat{n}\cdot\mathbf{x}'}{c}\right) d^3x'.
\end{equation}

For a monochromatic source oscillating at frequency $\Omega = 2\omega$
(the dominant frequency from $B^2 \propto \sin^2(\omega t)$):
\begin{equation}
S(\mathbf{x}', t) = \tilde{S}(\mathbf{x}') e^{-i\Omega t} + \text{c.c.},
\end{equation}
the far field becomes:
\begin{equation}\label{eq:far-field}
\psi_1(\mathbf{x}, t) = -\frac{e^{i(\Omega r/c - \Omega t)}}{4\pi r}
  \int \tilde{S}(\mathbf{x}') e^{-i\Omega\hat{n}\cdot\mathbf{x}'/c}\,d^3x'
  + \text{c.c.}
\end{equation}

\subsection{Long-Wavelength Limit}

The radiation wavelength is:
\begin{equation}
\lambda_\psi = \frac{2\pi c}{\Omega} = \frac{c}{2\nu}
  = \frac{3 \times 10^8}{2 \times 10^9} = 0.15\,\text{m} = 15\,\text{cm}.
\end{equation}

Since the source size $L_{\rm source} \sim 14$\,cm is comparable to
$\lambda_\psi$, we are \emph{not} in the extreme long-wavelength limit.
We must keep the phase factor $e^{-i\Omega\hat{n}\cdot\mathbf{x}'/c}$.

However, for an order-of-magnitude estimate, the integral in
Eq.~\eqref{eq:far-field} is:
\begin{equation}
\mathcal{I}(\hat{n}) = \int \tilde{S}(\mathbf{x}')
  e^{-i\Omega\hat{n}\cdot\mathbf{x}'/c}\,d^3x'.
\end{equation}

%=============================================================================
\section{$\psi$-Wave Luminosity}\label{sec:luminosity}
%=============================================================================

\subsection{Energy Flux in $\psi$-Waves}

The $\psi$-field carries energy density and flux. From the DFD field equation
$\Box\psi = -(8\pi G/c^4)\mathcal{T}$, the kinetic Lagrangian density is:
\begin{equation}
\mathcal{L}_\psi = -\frac{c^4}{16\pi G}\,\partial_\mu\psi\,\partial^\mu\psi,
\end{equation}
giving energy density and Poynting-like flux:
\begin{equation}
u_\psi = \frac{c^4}{16\pi G}\left(\frac{\dot{\psi}^2}{c^2}
  + |\nabla\psi|^2\right), \qquad
\mathbf{S}_\psi = -\frac{c^4}{8\pi G}\,\dot{\psi}\,\nabla\psi.
\end{equation}

Note: the factor $c^4/(16\pi G) = c^4/(16\pi G)$ has dimensions of
$[\text{J/m}^3]$ when $\psi$ is dimensionless and spatial derivatives
have dimensions $[\text{m}^{-1}]$:
\begin{equation}
\frac{c^4}{16\pi G} = \frac{(3\times10^8)^4}{16\pi \times 6.674\times10^{-11}}
  = \frac{8.1\times10^{33}}{3.35\times10^{-9}} = 2.42 \times 10^{42}\,
  \text{J/m}^3.
\end{equation}

This is the ``Planck energy density'' --- an enormous prefactor, but it
multiplies a tiny $\psi_1$.

\subsection{Radiated Power from the Far Field}

For a plane wave $\psi_1 = A\,e^{i(kr - \Omega t)}/r$, the time-averaged
energy flux is:
\begin{equation}
\langle S_\psi \rangle = \frac{c^3}{8\pi G}\,\Omega^2\,\frac{|A|^2}{r^2}.
\end{equation}

The total radiated power is:
\begin{equation}
P_\psi = \oint \langle S_\psi \rangle\,r^2\,d\Omega_{\hat{n}}
  = \frac{c^3\Omega^2}{8\pi G} \oint |A(\hat{n})|^2\,d\Omega_{\hat{n}}.
\end{equation}

From Eq.~\eqref{eq:far-field}, $A(\hat{n}) = \mathcal{I}(\hat{n})/(4\pi)$, so:
\begin{equation}\label{eq:total-power}
\boxed{P_\psi = \frac{c^3\Omega^2}{128\pi^3 G}
  \oint |\mathcal{I}(\hat{n})|^2\,d\Omega_{\hat{n}}}
\end{equation}

\subsection{Evaluating the Source Integral}

The spatial Fourier amplitude of the source:
\begin{equation}
\tilde{S}(\mathbf{x}') = \frac{4\pi G}{c^4}\,\kappa\,
  U_{\rm EM}(\mathbf{x}')\,\tilde{f}_\Omega\,
  \Theta[\eta(\mathbf{x}') - \etac],
\end{equation}
where $\tilde{f}_\Omega$ is the Fourier coefficient of $f^2(t)$ at
frequency $\Omega = 2\omega$. For $f(t) = |\sin(\omega t)|$:
\begin{equation}
f^2(t) = \sin^2(\omega t) = \frac{1}{2}(1 - \cos(2\omega t)),
\end{equation}
so $\tilde{f}_\Omega = 1/2$ at $\Omega = 2\omega$.

Therefore:
\begin{equation}
\mathcal{I}(\hat{n}) = \frac{2\pi G\kappa}{c^4} \int_{\eta > \etac}
  U_{\rm EM}(\mathbf{x}')\,e^{-i\Omega\hat{n}\cdot\mathbf{x}'/c}\,d^3x'.
\end{equation}

\paragraph{Order-of-magnitude estimate.}
Approximating the integral by a characteristic value:
\begin{equation}
|\mathcal{I}| \sim \frac{2\pi G\kappa}{c^4}\,\bar{U}_{\rm EM}\,V_{\rm at},
\end{equation}
where $\bar{U}_{\rm EM} = B_{\max}^2/(2\mu_0)$ is the peak EM energy
density and $V_{\rm at}$ is the above-threshold volume.

\paragraph{Numerical values.}
\begin{align}
\bar{U}_{\rm EM} &= \frac{(10)^2}{2 \times 4\pi \times 10^{-7}}
  = \frac{100}{2.51 \times 10^{-6}} = 3.98 \times 10^{7}\,\text{J/m}^3, \\
V_{\rm at} &\sim \pi r_{\rm sol}^2 \times (z_{\rm open} + z_{\rm shield})
  \approx \pi (0.02)^2 \times 0.138 = 1.73 \times 10^{-4}\,\text{m}^3, \\
|\mathcal{I}| &\sim \frac{2\pi \times 6.674 \times 10^{-11} \times \kappa}
  {(3 \times 10^8)^4} \times 3.98 \times 10^7 \times 1.73 \times 10^{-4} \\
  &= \frac{4.19 \times 10^{-10}\,\kappa}{8.1 \times 10^{33}}
  \times 6.89 \times 10^3 \\
  &= \frac{2.89 \times 10^{-6}\,\kappa}{8.1 \times 10^{33}}
  = 3.56 \times 10^{-40}\,\kappa\,\,[\text{m}^{-1}].
\end{align}

\subsection{Total Radiated Power}

Substituting into Eq.~\eqref{eq:total-power} with $\Omega = 2\pi \times 2
\times 10^9 = 4\pi \times 10^9$\,rad/s:
\begin{align}
P_\psi &= \frac{c^3\Omega^2}{128\pi^3 G}\,4\pi\,|\mathcal{I}|^2
  \qquad\text{(isotropic approximation for total)} \\
  &= \frac{c^3\Omega^2}{32\pi^2 G}\,|\mathcal{I}|^2.
\end{align}

Numerically:
\begin{align}
\Omega^2 &= (4\pi \times 10^9)^2 = 1.58 \times 10^{20}\,\text{rad}^2/\text{s}^2, \\
\frac{c^3\Omega^2}{32\pi^2 G} &= \frac{2.7 \times 10^{25} \times 1.58 \times 10^{20}}
  {32 \times 9.87 \times 6.674 \times 10^{-11}} \\
  &= \frac{4.27 \times 10^{45}}{2.11 \times 10^{-8}}
  = 2.02 \times 10^{53}\,[\text{J}\cdot\text{m}^2/\text{s}], \\
P_\psi &= 2.02 \times 10^{53} \times (3.56 \times 10^{-40})^2\,\kappa^2 \\
  &= 2.02 \times 10^{53} \times 1.27 \times 10^{-79}\,\kappa^2 \\
  &= 2.56 \times 10^{-26}\,\kappa^2\,\text{W}.
\end{align}

\begin{equation}\label{eq:power-result}
\boxed{P_\psi \approx 2.6 \times 10^{-26}\,\kappa^2\,\text{W}}
\end{equation}

For $\kappa = 1$, the radiated $\psi$-wave power is $\sim 3 \times 10^{-26}$\,W.

%=============================================================================
\section{Momentum Flux and Asymmetry}\label{sec:momentum}
%=============================================================================

\subsection{$\psi$-Field Momentum Density}

The $\psi$-field carries momentum density:
\begin{equation}\label{eq:momentum-density}
\mathbf{g}_\psi = \frac{1}{c^2}\mathbf{S}_\psi
  = -\frac{c^2}{8\pi G}\,\dot{\psi}_1\,\nabla\psi_1.
\end{equation}

For a radiation field propagating in direction $\hat{n}$:
\begin{equation}
\mathbf{g}_\psi = \frac{u_\psi}{c}\,\hat{n},
\end{equation}
as for any massless wave.

\subsection{Asymmetric Radiation Pattern}

The key physics: the above-threshold source volume is \emph{not}
centrosymmetric. More source volume lies on the open side ($z > 0$)
than the shielded side ($z < 0$).

Define the momentum radiated per unit time in direction $\hat{n}$:
\begin{equation}
\frac{dF}{d\Omega_{\hat{n}}} = \frac{r^2}{c}\langle S_\psi(\hat{n})\rangle
  = \frac{c^2\Omega^2}{8\pi G}\,\frac{|\mathcal{I}(\hat{n})|^2}{(4\pi)^2}.
\end{equation}

The net recoil force on the source is the integral of momentum flux
over all solid angles:
\begin{equation}\label{eq:recoil}
\mathbf{F}_{\rm recoil} = -\frac{1}{c} \oint \langle \mathbf{S}_\psi \rangle\,
  r^2\,d\Omega = -\frac{1}{c} \oint \hat{n}\,\langle S_\psi \rangle\,
  r^2\,d\Omega.
\end{equation}

\subsection{Asymmetry Factor}

The source integral $\mathcal{I}(\hat{n})$ depends on direction through
the phase factor $e^{-i\Omega\hat{n}\cdot\mathbf{x}'/c}$. For an axially
symmetric source along $\hat{z}$:
\begin{equation}
\mathcal{I}(\theta) = \frac{2\pi G\kappa}{c^4} \int_{\eta > \etac}
  U_{\rm EM}(r', z')\,e^{-i k z'\cos\theta}\,r'\,dr'\,dz',
\end{equation}
where $k = \Omega/c$ and $\theta$ is measured from $+\hat{z}$ (open end).

The crucial observation: for a source concentrated asymmetrically
(center of EM energy displaced toward $+\hat{z}$), the radiation pattern
$|\mathcal{I}(\theta)|^2$ is peaked in the forward ($+\hat{z}$) direction
due to constructive interference.

\paragraph{Dipole approximation for the asymmetry.}
Expand to first order in $kz'$:
\begin{equation}
\mathcal{I}(\theta) \approx \frac{2\pi G\kappa}{c^4}\left[
  \mathcal{M}_0 - ik\cos\theta\,\mathcal{M}_1\right],
\end{equation}
where:
\begin{align}
\mathcal{M}_0 &= \int_{\eta > \etac} U_{\rm EM}\,d^3x'
  \equiv \bar{U}\,V_{\rm at} \qquad\text{(monopole)}, \\
\mathcal{M}_1 &= \int_{\eta > \etac} z'\,U_{\rm EM}\,d^3x'
  \equiv \bar{U}\,V_{\rm at}\,\bar{z} \qquad\text{(dipole)},
\end{align}
with $\bar{z}$ the energy-weighted centroid displacement from the
geometric center.

The radiation power pattern is:
\begin{equation}
|\mathcal{I}(\theta)|^2 \approx \left(\frac{2\pi G\kappa}{c^4}\right)^2
  \left[\mathcal{M}_0^2 + k^2\cos^2\theta\,\mathcal{M}_1^2
  + 2k\cos\theta\,\mathcal{M}_0\,\text{Im}(\mathcal{M}_1)\right].
\end{equation}

Since $U_{\rm EM}$ is real, $\text{Im}(\mathcal{M}_1) = 0$. The
monopole--dipole cross term that would give net momentum vanishes
at this order!

\paragraph{Higher-order contribution.}
Net momentum transfer requires interference between modes of adjacent
$\ell$ (monopole-dipole or dipole-quadrupole). Since $\mathcal{M}_0$
is real and $\mathcal{M}_1$ is real, the leading asymmetric contribution
comes from the full phase factor. We write:
\begin{equation}
|\mathcal{I}(\theta)|^2 - |\mathcal{I}(\pi - \theta)|^2
  = 4\,\text{Re}\!\left[\mathcal{I}_0^*\,\mathcal{I}_1\right]\cos\theta
  + \ldots
\end{equation}
where $\mathcal{I}_0$ and $\mathcal{I}_1$ are monopole and dipole terms
with full complex structure.

For a source that is not long-wavelength ($kL \sim 1$), the asymmetry
in radiated momentum is:
\begin{equation}
\mathcal{A} \equiv \frac{|\text{momentum radiated}|}{P_\psi/c}
  \sim \frac{k\bar{z}}{1 + (k\bar{z})^2} \sim k\bar{z}
  \quad\text{for}\quad k\bar{z} \lesssim 1.
\end{equation}

\subsection{Numerical Estimate of the Asymmetry Factor}

The energy-weighted centroid:
\begin{equation}
\bar{z} = \frac{\int z\,U_{\rm EM}\,\Theta(\eta - \etac)\,d^3x}
  {\int U_{\rm EM}\,\Theta(\eta - \etac)\,d^3x}.
\end{equation}

For the asymmetric geometry, the above-threshold volume extends from
$z = -1.6$\,cm (shielded) to $z = +12.2$\,cm (open). If the energy
density is roughly uniform in the core and decays outside:
\begin{equation}
\bar{z} \approx \frac{12.2 - 1.6}{2}\,\text{cm} \approx 5.3\,\text{cm}
  = 0.053\,\text{m}.
\end{equation}

The wavevector:
\begin{equation}
k = \frac{\Omega}{c} = \frac{4\pi \times 10^9}{3 \times 10^8}
  = 41.9\,\text{m}^{-1}.
\end{equation}

Therefore:
\begin{equation}
k\bar{z} = 41.9 \times 0.053 = 2.22.
\end{equation}

This is order unity, meaning we are in the regime of significant
asymmetry. The asymmetry factor is:
\begin{equation}
\mathcal{A} \sim \frac{k\bar{z}}{1 + (k\bar{z})^2}
  = \frac{2.22}{1 + 4.93} = 0.37.
\end{equation}

A more careful numerical integration of the actual field profile would
be needed for precision, but $\mathcal{A} \sim 0.3$--$0.5$ is robust.

%=============================================================================
\section{Mode II Recoil Thrust}\label{sec:thrust}
%=============================================================================

\subsection{Thrust Formula}

The net recoil force is:
\begin{equation}
F_{\rm recoil} = \frac{\mathcal{A}\,P_\psi}{c},
\end{equation}
where $\mathcal{A}$ is the asymmetry factor and $P_\psi$ is the total
radiated $\psi$-wave power.

\subsection{Numerical Result}

\begin{align}
F_{\rm recoil} &= \frac{0.37 \times 2.6 \times 10^{-26}\,\kappa^2}
  {3 \times 10^8} \\
  &= 3.2 \times 10^{-35}\,\kappa^2\,\text{N}.
\end{align}

For $\kappa = 1$:
\begin{equation}\label{eq:thrust-result}
\boxed{F_{\rm Mode\,II} \approx 3 \times 10^{-35}\,\text{N}}
\end{equation}

\subsection{Scaling Law}

Collecting all factors, the Mode II thrust scales as:
\begin{equation}\label{eq:scaling}
\boxed{F_{\rm Mode\,II} = \frac{\kappa^2 G \Omega^2}{32\pi^2 c^5}
  \left(\bar{U}_{\rm EM}\,V_{\rm at}\right)^2 \mathcal{A}}
\end{equation}

This confirms \textbf{option (c)} from the problem statement: the source
amplitude involves $\kappa$ (not $\lambda$), but the radiation formula
inherently involves $G/c^5$ because $\psi$-waves are gravitational-type
scalar waves.

\paragraph{Explicit parameter dependence:}
\begin{equation}
F \propto \kappa^2 \times \frac{G}{c^5} \times (U_{\rm EM}\,V_{\rm at})^2
  \times \omega^2 \times \mathcal{A}.
\end{equation}

The factor $G/c^5 = 6.674 \times 10^{-11}/(2.43 \times 10^{42})
= 2.75 \times 10^{-53}$\,s/J is the fundamental reason why scalar wave
radiation from laboratory sources is negligibly small.

%=============================================================================
\section{Comparison of Thrust Channels}\label{sec:comparison}
%=============================================================================

\begin{table}[h]
\centering
\caption{Mode II thrust through different channels.}\label{tab:channels}
\begin{tabular}{lccc}
\toprule
Channel & Coupling & Suppression & Thrust (N) \\
\midrule
$\lambda$ (Appendix R) & $|\lambda-1| \lesssim 3\times10^{-5}$ & $(|\lambda-1|)^2 \times G/c^5$ & $\lesssim 10^{-44}$ \\
$\kappa$ (dual-sector) & $|\kappa| \lesssim 1$ & $\kappa^2 \times G/c^5$ & $\sim 10^{-35}$ \\
\bottomrule
\end{tabular}
\end{table}

Both channels produce negligible thrust because the factor $G/c^5$
suppresses any laboratory-scale scalar wave emission.

%=============================================================================
\section{Can the $G/c^5$ Suppression Be Circumvented?}\label{sec:circumvent}
%=============================================================================

\subsection{The Fundamental Issue}

The luminosity formula $P \propto G/c^5 \times (\text{source})^2$ arises
because $\psi$-waves are gravitational-type scalar waves. Their kinetic
term is $\mathcal{L} \propto (c^4/G)\partial\psi\,\partial\psi$, giving
the canonically normalized field $\phi = (c^2/\sqrt{G})\psi$. Any source
couples to $\psi$ through $G$, and the radiation formula inherits $G/c^5$.

This is identical to the reason gravitational waves from laboratory
sources are undetectable: the $G/c^5$ factor is $\sim 10^{-53}$\,s/J.

\subsection{Could There Be a Non-Gravitational Channel?}

For Mode II thrust to be macroscopic, one would need a $\psi$-sourcing
mechanism that does \emph{not} go through $G$. This would require:
\begin{enumerate}
\item A direct coupling $\psi \leftrightarrow$ EM with strength much
  larger than $G/c^4$.
\item Such a coupling would need to be consistent with all existing
  constraints (cavity stability, Shapiro delay, etc.).
\end{enumerate}

From Appendix R, the $\lambda$ channel provides exactly such a direct
coupling --- but it is bounded by $|\lambda - 1| \lesssim 3 \times 10^{-5}$,
and the natural value is $\lambda = 1$ (no coupling).

The $\kappa$ channel modifies how EM responds to $\psi$ (not how EM
sources $\psi$). The sourcing still goes through the gravitational
$\psi$-field equation with its $G/c^4$ coupling.

\subsection{Verdict}

Within the DFD framework as currently formulated, all $\psi$-wave
radiation from EM sources is gravitationally suppressed. There is no
known mechanism to produce macroscopic Mode II thrust.

%=============================================================================
\section{What If $\neff$ Modification Does Not Radiate?}\label{sec:mode1}
%=============================================================================

An alternative interpretation: the $\neff$ modification of
Eq.~\eqref{eq:neff} is a \emph{local} change to the optical metric,
not a propagating $\psi$-wave. In this case:

\begin{itemize}
\item The pulsed $\neff$ modification creates a time-varying local metric
  within the above-threshold volume.
\item This does \emph{not} radiate scalar waves (it is a driven,
  non-propagating perturbation).
\item The recoil mechanism would be Mode I (direct EM momentum asymmetry
  from the modified metric), not Mode II (radiated $\psi$-wave recoil).
\end{itemize}

Mode I analysis is the subject of a separate calculation. The key
difference: Mode I forces scale with the EM energy and the local
$\psi$-gradient, without the $G/c^5$ suppression.

%=============================================================================
\section{Summary and Conclusions}\label{sec:summary}
%=============================================================================

\begin{tcolorbox}[colback=blue!5!white, colframe=blue!50!black,
  title=Mode II Thrust: Key Results]
\begin{enumerate}
\item \textbf{Source channel:} The $\kappa$ parameter (dual-sector
  extension, Appendix R \S R.6) provides the unsuppressed EM--$\psi$
  coupling. The $\lambda$ channel is bounded to $|\lambda-1| < 3\times10^{-5}$.

\item \textbf{Radiation formula:} $\psi$-wave luminosity is:
\begin{equation*}
P_\psi = \frac{\kappa^2 G \Omega^2}{32\pi^2 c^5}
  (\bar{U}_{\rm EM}\,V_{\rm at})^2 \times \mathcal{G}(\theta),
\end{equation*}
where $\mathcal{G}$ is a geometric factor of order $4\pi$.

\item \textbf{Numerical result:} For the 10\,T asymmetric solenoid at
  1\,GHz in UHV:
\begin{equation*}
P_\psi \approx 3 \times 10^{-26}\,\kappa^2\,\text{W}, \qquad
F_{\rm recoil} \approx 3 \times 10^{-35}\,\kappa^2\,\text{N}.
\end{equation*}

\item \textbf{Scaling:} $F \propto \kappa^2 \times (G/c^5) \times
  (U_{\rm EM} V)^2 \omega^2 \times \mathcal{A}$.

\item \textbf{Verdict:} The $G/c^5 \approx 3 \times 10^{-53}$\,s/J
  prefactor makes Mode II thrust gravitationally negligible. This is
  the same suppression that makes laboratory GW emission undetectable.

\item \textbf{Implication for drive concept:} Mode II (scalar wave
  recoil) is not viable. Any $\psi$-bubble drive must operate through
  Mode I (direct metric modification) or a fundamentally different
  mechanism.
\end{enumerate}
\end{tcolorbox}

\begin{tcolorbox}[colback=red!5!white, colframe=red!50!black,
  title=Critical Finding]
\textbf{The $G/c^5$ barrier is fundamental}, not an artifact of weak
coupling. It arises from the normalization of the $\psi$-field kinetic
term. The only way to circumvent it would be if $\psi$ had an additional,
non-gravitational kinetic term --- which is not part of the current DFD
framework.

The correct path for the $\psi$-bubble drive is Mode I: direct, local
modification of the effective metric through $\neff$, which does not
require $\psi$-wave radiation and avoids the $G/c^5$ suppression entirely.
\end{tcolorbox}

%=============================================================================
\appendix
\section{Dimensional Analysis Cross-Check}\label{app:dimensions}
%=============================================================================

\paragraph{Source term.}
$[S] = [G/c^4]\times[U_{\rm EM}]\times[\kappa] =
(\text{m}^3/(\text{kg}\cdot\text{s}^2))/(\text{m}^4/\text{s}^4)
\times (\text{J/m}^3) = \text{m}^{-2}$. $\checkmark$ (matches $[\Box\psi]$)

\paragraph{Source integral.}
$[\mathcal{I}] = [S]\times[\text{m}^3] = \text{m}^{-2}\times\text{m}^3 = \text{m}$.
But from Eq.~\eqref{eq:far-field}, $A = \mathcal{I}/(4\pi)$ and $\psi_1 = A\,e^{ikr}/r$,
so $[\mathcal{I}] = [\psi_1]\times[\text{m}] = \text{m}$ (dimensionless $\times$ m). $\checkmark$

\paragraph{Power.}
$[c^3\Omega^2/(G)] \times [\mathcal{I}]^2 =
(\text{m}^3/\text{s}^3)\times(\text{s}^{-2})/(\text{m}^3/(\text{kg}\cdot\text{s}^2))
\times \text{m}^2 = \text{kg}\cdot\text{m}^2/\text{s}^3 = \text{W}$. $\checkmark$

\paragraph{Thrust.}
$[P_\psi/c] = \text{W}/(\text{m/s}) = \text{N}$. $\checkmark$

%=============================================================================
\section{Parameter Table}\label{app:parameters}
%=============================================================================

\begin{table}[h]
\centering
\caption{Input parameters for the Mode II calculation.}
\begin{tabular}{lcc}
\toprule
Parameter & Symbol & Value \\
\midrule
Peak magnetic field & $B_{\max}$ & 10\,T \\
Vacuum gas density & $\rho$ & $10^{-8}$\,kg/m$^3$ \\
Shielded decay length & $L_s$ & 1\,cm \\
Open decay length & $L_o$ & 10\,cm \\
Above-threshold (open) & $z_{\rm open}$ & 12.2\,cm \\
Above-threshold (shielded) & $z_{\rm shield}$ & 1.6\,cm \\
Pulse frequency & $\nu$ & 1\,GHz \\
Angular frequency & $\omega$ & $2\pi \times 10^9$\,rad/s \\
$\psi$-wave frequency & $\Omega = 2\omega$ & $4\pi \times 10^9$\,rad/s \\
DFD threshold & $\etac = \alpha/4$ & $1.82 \times 10^{-3}$ \\
$\kappa$ coupling & $|\kappa|$ & $\lesssim 1$ \\
$\lambda$ coupling & $|\lambda - 1|$ & $\lesssim 3\times10^{-5}$ \\
Asymmetry ratio & $\mathcal{R}_{\rm geom}$ & 7.6 \\
Asymmetry factor & $\mathcal{A}$ & $\sim 0.37$ \\
Peak EM energy density & $\bar{U}_{\rm EM}$ & $4.0 \times 10^7$\,J/m$^3$ \\
Above-threshold volume & $V_{\rm at}$ & $1.7 \times 10^{-4}$\,m$^3$ \\
\midrule
\textbf{Result: Radiated power} & $P_\psi$ & $2.6 \times 10^{-26}\kappa^2$\,W \\
\textbf{Result: Recoil thrust} & $F_{\rm Mode\,II}$ & $3.2 \times 10^{-35}\kappa^2$\,N \\
\bottomrule
\end{tabular}
\end{table}

\end{document}
